\documentclass[11pt]{article}
\usepackage[utf8]{inputenc}
\usepackage[T1]{fontenc}
\usepackage[margin=1in]{geometry}
\usepackage{hyperref}
\usepackage{enumitem}
\hypersetup{colorlinks=true,linkcolor=black,urlcolor=blue,citecolor=blue}
\title{Interactive Word Wall Build}
\author{Piter Garcia}
\date{March 20, 2026}
\begin{document}
\maketitle
\section*{Packet Use}
This printable lesson packet is generated from the paired Markdown guide. The Markdown version remains the clickable, neurodivergent-friendly working copy; this PDF carries the same lesson substance in a formal handout format. Evidence claims in this packet are based on transcript-derived lesson notes and the cleaned transcript files in this workspace.

\section*{Quick Scan}
\begin{itemize}[leftmargin=*]
\item \textbf{Grade:} Grade 1
\item \textbf{Grade Hub:} Notes [../grade-1-Notes.md]
\item \textbf{Schedule Reference:} Spring2026\_TeachingSchedule\_Derek.md [../../school/Spring2026\_TeachingSchedule\_Derek.md]
\item \textbf{Workbook Sheet:} \texttt{1st} in \texttt{Teaching\_Placement-IGNITE Curriculum\_ GCSD25-26.xlsx}
\item \textbf{Lesson Focus:} Students build an interactive vocabulary display that links words, visuals, and classroom-made media through QR or video-supported artifacts.
\item \textbf{CS Alignment Strength:} CS-adjacent digital communication and information representation
\item \textbf{LaTeX Source:} interactive-word-wall-build.tex [./interactive-word-wall-build.tex]
\item \textbf{Bib File:} interactive-word-wall-build.bib [./interactive-word-wall-build.bib]
\item \textbf{Artifacts Folder:} artifacts/ [./artifacts/]
\end{itemize}

\section*{Lesson Identity}
\begin{itemize}[leftmargin=*]
\item \textbf{Likely Class Blocks:} Day 3 -- 12:25-1:15 -- Gargana, Day 3 -- 2:15-3:05 -- Polfleit, Day 5 -- 12:25-1:15 -- Montgomery
\item \textbf{Main Lesson Family:} Interactive Word Wall Build
\item \textbf{Primary Evidence Base:} Transcript-derived notes plus source transcripts from this teaching-placement workspace.
\item \textbf{Primary External Source Type:} Official curriculum/provider materials.
\end{itemize}

\section*{What We Are Teaching}
\begin{itemize}[leftmargin=*]
\item representing information with text, image, and object together
\item using a predictable build sequence to assemble a shared classroom artifact
\item linking a vocabulary item to a digital or recorded explanation
\item participating in collaborative classroom technology routines
\end{itemize}

\section*{CS Standards Alignment}
\begin{itemize}[leftmargin=*]
\item \textbf{1A-DA-06}: Collect and present the same data in various visual formats. Why it fits here: Students represent one idea across word cards, images, QR codes, and built pieces.
\item \textbf{1A-AP-08}: Model daily processes by creating and following algorithms. Why it fits here: Students follow a repeatable sequence for selecting, building, and attaching their component.
\item \textbf{1A-IC-17}: Work respectfully and responsibly with others when using shared technology. Why it fits here: The lesson depends on collaboration, turn-taking, and shared artifact care.
\end{itemize}

\section*{NYS / Local Curriculum Alignment}
\begin{itemize}[leftmargin=*]
\item Aligned to the local IGNITE first-grade sheet through literacy-tech integration and classroom digital creation.
\item Supports New York-facing communication and media fluency by requiring students to connect a concept to a digital or scannable representation.
\item This is a lighter-CS lesson, but it still counts as computing instruction because students organize information structures and use technology-supported representation intentionally.
\end{itemize}

\section*{Lesson Content and Concepts}
\begin{itemize}[leftmargin=*]
\item representation
\item media linking
\item sequence
\item shared artifact
\item information organization
\end{itemize}

\section*{Evidence / Proof of Instruction}
\begin{itemize}[leftmargin=*]
\item \textbf{Transcript-derived note:} 260211 Grade 1 Gargana Word Wall [../../transcript-derived-lessons/260211-grade-1-gargana-word-wall.md] The lesson note captures the Google Vids / QR word wall setup as an instructional design choice, not a casual craft activity.
\item \textbf{Transcript-derived note:} 260212 Grade 1 Gargana Lego Wall [../../transcript-derived-lessons/260212-grade-1-gargana-lego-wall.md] The second note shows that the word wall idea persisted as a structured project across sessions.
\item \textbf{Source transcript:} 260211 Grade 1 Gargana Word Wall [../../transcripts/260211-grade-1-gargana-word-wall.md] The transcript mentions objectives, LEGO-linked word wall pieces, and teacher concern about giving students explicit directions and outcomes.
\item \textbf{Likely student proof:} Built vocabulary pieces, QR-linked media, recorded explanations, and the finished wall display all serve as evidence of teaching and learning.
\end{itemize}

\section*{Assessment Evidence}
\begin{itemize}[leftmargin=*]
\item Student can explain what their piece represents.
\item Student can match a word to a visual or recorded explanation.
\item Student participates in the shared build without breaking the collaboration routine.
\end{itemize}

\section*{UDL and Inclusion Supports}
\subsection*{Multiple Representation}
\begin{itemize}[leftmargin=*]
\item Pair every word with a picture, color cue, and concrete object.
\item Use teacher-modeled examples before students create their own pieces.
\item Keep instructions anchored to one visible exemplar at a time.
\end{itemize}

\subsection*{Multiple Action / Expression}
\begin{itemize}[leftmargin=*]
\item Allow drawing, speaking, pointing, or building as valid ways to show understanding.
\item Use partner roles: builder, holder, explainer, scanner.
\item Offer pre-cut or partially prepared materials for students who fatigue quickly.
\end{itemize}

\subsection*{Multiple Engagement}
\begin{itemize}[leftmargin=*]
\item Let students choose which word or concept they represent.
\item Use scanning, recording, or showing as motivating interactive elements.
\item Celebrate each added piece so the wall becomes a visible cumulative success marker.
\end{itemize}

\section*{How We Are Already Making It Inclusive}
\begin{itemize}[leftmargin=*]
\item Transcript evidence suggests the lesson already used strong teacher objectives and concrete artifact goals.
\item The build format supports students who need movement and tactile entry points.
\item Using multiple media forms already broadens access beyond text-only vocabulary work.
\end{itemize}

\section*{How to Improve Inclusion Next Time}
\begin{itemize}[leftmargin=*]
\item Add a one-page visual build checklist at each table.
\item Provide bilingual vocabulary cards and recorded directions.
\item Offer a no-QR fallback path where students can point to a recorded explanation instead of scanning independently.
\item Save sample student pieces in \texttt{artifacts/} to strengthen proof of instruction.
\end{itemize}

\section*{Materials and Setup}
\begin{itemize}[leftmargin=*]
\item Word cards
\item QR links or short recorded videos
\item Maker materials / LEGO or build pieces
\item Display space
\item Student or teacher scanning device
\end{itemize}

\section*{Teacher Moves / Script / Routines}
\begin{itemize}[leftmargin=*]
\item Start with one finished example and reverse-engineer how it was made.
\item Assign one clear role before materials are distributed.
\item Pause the room to compare two strong representation choices.
\item Close by having students explain how the wall helps another learner understand the word.
\end{itemize}

\section*{Common Errors and Debugging}
\begin{itemize}[leftmargin=*]
\item Students may focus on decoration instead of meaning.
\item QR or video links may fail without a backup display plan.
\item Shared materials can create conflict if roles are not explicit.
\end{itemize}

\section*{Extensions and Simplifications}
\begin{itemize}[leftmargin=*]
\item Extension: add a short student-recorded explanation to each entry.
\item Simplification: limit each team to one word, one image, and one explanation.
\item Extension: sort wall entries into categories and ask students to explain the pattern.
\end{itemize}

\section*{Related Local Materials}
\begin{itemize}[leftmargin=*]
\item No direct local lesson packet linked yet.
\end{itemize}

\section*{Artifacts Checklist}
\begin{itemize}[leftmargin=*]
\item Compiled lesson PDF in \texttt{artifacts/}
\item At least one screenshot, student example, or setup photo
\item One short assessment or observation record
\item Updated standards/source bibliography once the \texttt{.bib} file is revised
\item Any teacher reflection or revision notes after reteaching
\end{itemize}

\section*{Selected Official Sources}
Selected official sources that informed this packet are listed below.
ocite{*}

\bibliographystyle{plain}
\bibliography{interactive-word-wall-build}
\end{document}
